\providecommand{\stelladifficolta}[1]{%
  \tikz[baseline=-0.60ex, x=0.48em, y=0.48em]{%
    \path[draw=black, fill=#1, line width=0.3pt, line join=round]
      (90:1) -- (54:0.4) -- (18:1) -- (-18:0.4) -- (-54:1)
      -- (-90:0.4) -- (-126:1) -- (-162:0.4) -- (-198:1) -- (-234:0.4)
      -- cycle;%
  }%
}
\providecommand{\stellavalutata}[2]{%
  \ifnum#1<#2\stelladifficolta{white}\else\stelladifficolta{black}\fi%
}
\providecommand{\valutazionestelle}[1]{%
  \begingroup
  \stellavalutata{#1}{1}%
  \stellavalutata{#1}{2}%
  \stellavalutata{#1}{3}%
  \stellavalutata{#1}{4}%
  \stellavalutata{#1}{5}%
  \endgroup}
\providecommand{\difficoltaesercizio}[1]{%
  \leavevmode\unskip\hfill\valutazionestelle{#1}\par\smallskip}
\providecommand{\rimandoteoria}[1]{%
  \par\smallskip\noindent{\footnotesize\emph{Riferimento negli appunti:} \cref{#1}.}\par}

\phantomsection
\section*{Scala di difficoltà degli esercizi}
\addcontentsline{toc}{section}{Scala di difficoltà degli esercizi}
\label{sec:scala-difficolta-esercizi}

La valutazione è soggettiva e riflette l'impegno richiesto e le difficoltà
incontrate durante lo svolgimento. Non rappresenta una misura oggettiva
della difficoltà né un'indicazione dell'importanza dell'esercizio
nel programma e può quindi differire da quella di altri studenti.

\begin{center}
\small
\renewcommand{\arraystretch}{1.45}
\begin{tabular}{@{}p{0.19\linewidth}p{0.18\linewidth}p{0.55\linewidth}@{}}
\hline
\textbf{Stelle} & \textbf{Livello} & \textbf{Descrizione}\\
\hline
\valutazionestelle{0} & Immediato &
Si risolve a prima vista, senza richiedere ragionamenti o procedimenti particolari.\\
\valutazionestelle{1} & Elementare &
Richiede l'applicazione diretta di conoscenze o procedimenti noti, senza particolari difficoltà.\\
\valutazionestelle{2} & Facile &
Richiede alcuni passaggi e un ragionamento autonomo, ma la strategia risolutiva rimane abbastanza evidente.\\
\valutazionestelle{3} & Intermedio &
Richiede di individuare una strategia non immediata e sviluppare un ragionamento articolato, eventualmente attraverso più tentativi.\\
\valutazionestelle{4} & Difficile &
Richiede un ragionamento impegnativo, una strategia poco evidente e diversi tentativi per giungere alla soluzione.\\
\valutazionestelle{5} & Complesso &
Richiede intuizioni particolarmente difficili e numerosi tentativi; alcuni passaggi possono risultare ostici anche dopo aver consultato la soluzione.\\
\hline
\end{tabular}
\end{center}

\clearpage
